Here $\Omega$ is the rotation-gauge field and $R$ its field strength (Fact \claimref{synthesis.imported-gtg-dictionary}(ii) and (iv)), so every step is a computation in Gauge Theory Gravity on one flat chart; the component readings use $\Omega = \tfrac12\omega^{ab}\gamma_a\wedge\gamma_b$ and $R(a\wedge b) = \tfrac12 R^{ab}(a, b)\,\gamma_a\wedge\gamma_b$. (i) is the transgression formula for a quadratic invariant pairing, here $\langle\cdot,\cdot\rangle_{\epsilon}$ on the bivectors, $\mathfrak{so}(1,3)$: the pairing is symmetric and invariant, $\langle Z\times X, Y\rangle_{\epsilon} + \langle X, Z\times Y\rangle_{\epsilon} = 0$, since $I_4$ commutes with bivectors. The exterior derivative of the derivative term of $\mathrm{CS}$ is the pairing of the derivative part of $R$ with itself, that of the cubic term is $\tfrac23\cdot 3$ times the pairing of the derivative part with $\Omega\times\Omega$, so $\nabla\wedge\mathrm{CS}$ is $\mathcal{E}_{R}$ up to the pairing of $\Omega\times\Omega$ with itself antisymmetrised over four arguments, which vanishes by invariance and the Jacobi identity (in components $d(\omega\,d\omega) = d\omega\wedge d\omega$, $d(\omega^3) = 3\,d\omega\wedge\omega^2$ and $\epsilon(\omega^2,\omega^2) = 0$). The script confirms $d\,\mathrm{CS} = \epsilon_{abcd}R^{ab}\wedge R^{cd}$ as an exact identity of forms on the Kasner connection and on a generic $\mathfrak{so}(1,3)$ connection. (ii) On the slice only $\mathrm{CS}(e_1\wedge e_2\wedge e_3)$ enters. Its derivative term vanishes there for two reasons, either sufficient: the field depends on $\tau$ alone, so $b\cdot\nabla\Omega(c) = 0$ for tangential $b$, $c$; and every nonzero value of $\Omega$ is a boost $B_i$, while $\langle B_i, B_j\rangle_{\epsilon} = 0$ ($I_4B_i$ is a multiple of the complementary plane). The cubic term is $2\langle\Omega(e_1), \Omega(e_2)\times\Omega(e_3)\rangle_{\epsilon}$ with $\Omega(e_i) = -\dot a_iB_i$ on background directions, $a_i = \tau^{p_i}$, and $B_2\times B_3 = -e_2\wedge e_3$, so it is $2\,p_1p_2p_3\,\tau^{\sum p_i - 3}\langle B_1, e_2\wedge e_3\rangle_{\epsilon} = 8\,p_1p_2p_3\,\tau^{-2}$, since $\langle B_1, e_2\wedge e_3\rangle_{\epsilon} = 4$ and $\sum p_i = 1$ (in components, the cubic term reduces to $\tfrac23\epsilon_{i0jk}\,\omega^{i0}\wedge\omega^{j}{}_{0}\wedge\omega^{0k}$ and its index rearrangements, each a multiple of $p_1p_2p_3\tau^{-3}e^1\wedge e^2\wedge e^3 = p_1p_2p_3\tau^{-2}d^3x$, with total coefficient $8$); the script computes the $(x,y,z)$ component of $\mathrm{CS}$ exactly and finds $8p_1p_2p_3\tau^{-2}$. Then $\int_{\tau>\varepsilon}\mathcal{E}_{R} = \int \mathrm{CS}\big|_{\text{far}} - \int_{\Sigma_\varepsilon}\mathrm{CS}$ by Stokes, and $-\tfrac3\Lambda\cdot 8p_1p_2p_3\varepsilon^{-2}$ integrates $-\tfrac{3}{\Lambda}\cdot 16 p_1p_2p_3\int_\varepsilon\tau^{-3}d\tau$, matching Proposition \claimref{dark-energy.frozen-point-datum-kasner-wall}(iii). (iii) A boundary term does not change the Euler-Lagrange equations in the interior. By (ii) the added term cancels the wall boundary contribution of the exact topological term identically on the Kasner approach, leaving the far-boundary value of $\mathrm{CS}$ and the two remaining terms of the expansion in \claimref{dark-energy.frozen-point-makes-stelle-west-vector}(2), the Einstein-Cartan term (twice the density $-2\,(\partial_a\cdot\mathcal{R}(a))\det(\hb)^{-1}$, which is $-8\Lambda\tau$ on the Einstein region; in components $2\int\epsilon_{abcd}e^ae^bR^{cd}$) and the cosmological term ($\tfrac{\Lambda}{3}$ times the density $24\det(\hb)^{-1} = 24\tau$; in components $\tfrac{\Lambda}{3}\int\epsilon_{abcd}e^ae^be^ce^d$), all finite as $\varepsilon\to 0$. With the corrections of Proposition \claimref{dark-energy.cosmological-constant-does-not-change-residue}, $\Omega(\hb(e_i)) = -(p_i/\tau + \tfrac{\Lambda}{2}(1-p_i)\tau + \dots)B_i$ (in components $\omega^{i0} = (p_i/\tau + \tfrac{\Lambda}{2}(1-p_i)\tau + \dots)e^i$) and the cubic term of $\mathrm{CS}$ on the slice acquires, beyond $8p_1p_2p_3\tau^{-2}$, a term of order $\Lambda\tau^0$ and higher: a finite constant, which is the only freedom in the subtraction.
